\par\noindent\textbf{(1)} 允许区局部波数k(x)=√[2m(E−V)]/ℏ，WKB波形 ψ≈k\textasciicircum{}\{−1/2\}[Ae\textasciicircum{}\{i∫kdx\}+Be\textasciicircum{}\{−i∫kdx\}]。禁阻区κ(x)=√[2m(V−E)]/ℏ，解为κ\textasciicircum{}\{−1/2\}exp(±∫κdx)。在k→0的转折点这些幅因子发散，必须用局域Airy方程连接而非硬延拓。
\par\noindent\textbf{(2)} 对简单线性转折点，Airy连接把向禁阻侧衰减的解映成允许区 2k\textasciicircum{}\{−1/2\}cos(∫\_\{x₁\}\textasciicircum{}xkdx−π/4)；另一端的衰减条件也贡献一个π/4。两端总Maslov相位为π/2，相比无转折点的单值周期条件产生半量子数。
\par\noindent\textbf{(3)} 令两侧允许区相位相容，得到 ∫\_\{x₁\}\textasciicircum{}\{x₂\}√[2m(E−V(x))]dx=πℏ(n+1/2)，n=0,1,2,…。等价经典作用量 ∮pdx=2πℏ(n+1/2)，其中每个简单转折点贡献Maslov指数1。
\par\noindent\textbf{(4)} 对V=mω²x²/2，转折点±a，a=√[2E/(mω)]；作用积分 ∫\_\{−a\}\textasciicircum{}a mω√(a²−x²)dx=πE/ω。量子化给E\_n=ℏω(n+1/2)，恰为精确谐振子谱。一般势需 |k'|≪k² 且转折点孤立；合并转折点或强量子隧穿时需一致近似。
